\section{An audit of video token reducers}
\label{sec:audit}

Proposition~\ref{prop:consecutive} predicts which reducers miss slow change. Here we test the prediction, and the converse question of what anchored rules miss, on planted events at matched budgets.

\paragraph{What is re-implemented.} No published system is run. Each method is reduced to the temporal comparison that decides whether a token is redundant, with the model-side statistic replaced by the pixel-space statistic that plays the same role (Appendix~\ref{app:reducers}). Seven methods reduce to four comparisons. \emph{Consecutive mean}: mean $|F_t - F_{t-1}|$ per patch below $\tau$, the pixel test of EVS \citep{evs2025} and a per-frame analogue of the normalised-pixel $L^1$ test of run-length tokenization \citep{choudhury2024rlt}; EVS's percentile cutoff is a threshold set by the budget, which is what calibration does. \emph{Consecutive cosine}: cosine similarity of patch vectors of $F_t$ and $F_{t-1}$ above $1 - \tau$, with mean-removed random projections standing in for a linear patch embedding: the feature test of TimeChat-Online \citep{yao2025timechatonline} and FrameFusion \citep{fu2024framefusion}, and the adjacent-frame part of DyCoke \citep{tao2025dycoke}. \emph{Segment static}: within fixed 16-frame segments, a patch position is static iff the mean pairwise cosine similarity of its patch vectors over the segment's frames is at least $1 - \tau$, and a static position is kept once per segment: the static-token test of PruneVid \citep{huang2025prunevid}, with fixed segments in place of its clustering. \emph{Window similarity}: within non-overlapping 8-frame windows, a whole frame is dropped iff its mean cosine similarity to the window's other frames is at least $1 - \tau$: the temporal reduction of LongVU \citep{shen2024longvu}, without its query-guided and spatial steps. Three reference rules complete the set: the consecutive mean with a full refresh every 60 frames, uniform frame sampling, and the anchored mean (mean $|F_t - R|$ against the last kept copy, the anchored rule with no certificate). The certified rules are the range rule, quotient + multi-scale at $\Delta_0 = 80$, and the same with the sequential layer \citep{nair2026certifiedskip}. What this tests is the comparison structure of each method, not its accuracy as a system: a learned feature can be more or less sensitive than a pixel statistic, but it cannot see a change against a frame it is not compared with.

\paragraph{Protocol.} Six kinds of event are planted into real fixed-camera footage (Appendix~\ref{app:events}): a small $5{\times}5$ object of contrast 40 for one frame (30 events); a tiny $3{\times}3$ object of contrast 48 for one frame (30); a fade of an $8{\times}8$ square to contrast 48 over 3~s (20); a slow fade of a $12{\times}12$ square to contrast 48 over 12~s (20); a $6{\times}6$ mover of contrast 40 at 2~px per frame for 12 frames (15); and a faint persistent $6{\times}6$ object of contrast 12 for 3~s (15). A one-frame or moving event is caught if its patch is kept while the object is present; a faint object if it is kept within 1~s; a fade if it is kept at or before the frame where the fade reaches contrast 32. Every rule's threshold is calibrated on the untouched clip to the drop rate of the certified quotient + multi-scale rule at $\Delta_0 = 80$, so all rules spend the same budget: every thresholded rule matches it to within 0.001. Three rules move in whole frames (uniform sampling, LongVU's frame rule, the periodic refresh) and are shown at their nearest setting, within 0.003 except the refresh on VIRAT, whose keyframes alone exceed the budget and which therefore gets 1.7--1.9 times the tokens of the others. Two clips (the CAVIAR lobby \citep{caviar2004} and VIRAT~S\_000200\_00 \citep{oh2011virat}) are run at their native rates (25 and 10~fps) and at 2~fps, the rate at which a video LLM is usually fed. The column $\rho$ is the rule's median false-fire rate on the fading patches before their fade starts.

\begin{table}[t]
\centering\scriptsize
\setlength{\tabcolsep}{3.5pt}
\resizebox{\linewidth}{!}{\begin{tabular}{lrrrrrrrrr}
\toprule
rule & drop & tokens/s & small 5$\times$5 & tiny 3$\times$3 & fade 3 s & fade 12 s & mover & faint 12 & $\rho$ median \\
\midrule
certified range & 0.976 & 263 & 0.07 & 0.13 & 0.15 & 0.25 & 0.08 & 0.13 & 0.000 \\
certified quotient+MS & 0.977 & 246 & 1.00 & 1.00 & 1.00 & 0.95 & 0.90 & 0.20 & 0.000 \\
certified +sequential & 0.977 & 251 & 1.00 & 1.00 & 1.00 & 0.95 & 0.91 & 1.00 & 0.001 \\
anchored mean (no certificate) & 0.977 & 246 & 0.27 & 0.10 & 1.00 & 1.00 & 0.41 & 0.13 & 0.000 \\
consecutive mean (EVS, RLT) & 0.977 & 246 & 0.10 & 0.03 & 0.05 & 0.15 & 0.16 & 0.07 & 0.000 \\
consecutive cosine (TimeChat, FrameFusion) & 0.977 & 246 & 0.53 & 0.40 & 0.20 & 0.25 & 0.42 & 0.13 & 0.000 \\
segment static (PruneVid) & \multicolumn{9}{l}{cannot reach the budget at this frame rate} \\
window similarity (LongVU) & 0.977 & 249 & 0.03 & 0.03 & 0.00 & 0.15 & 0.05 & 0.07 & 0.031 \\
consecutive + refresh every 60 & 0.977 & 246 & 0.07 & 0.03 & 0.75 & 1.00 & 0.04 & 0.33 & 0.016 \\
uniform frames & 0.978 & 238 & 0.07 & 0.07 & 1.00 & 1.00 & 0.02 & 0.40 & 0.022 \\
\bottomrule
\end{tabular}
}

\vspace{4pt}
\resizebox{\linewidth}{!}{\begin{tabular}{lrrrrrrrrr}
\toprule
rule & drop & tokens/s & small 5$\times$5 & tiny 3$\times$3 & fade 3 s & fade 12 s & mover & faint 12 & $\rho$ median \\
\midrule
certified range & 0.992 & 67 & 0.07 & 0.13 & 0.00 & 0.05 & 0.00 & 0.00 & 0.000 \\
certified quotient+MS & 0.992 & 68 & 1.00 & 1.00 & 1.00 & 0.90 & 0.93 & 0.00 & 0.000 \\
certified +sequential & 0.992 & 71 & 1.00 & 1.00 & 1.00 & 0.90 & 0.93 & 1.00 & 0.000 \\
anchored mean (no certificate) & 0.990 & 88 & 1.00 & 0.90 & 1.00 & 1.00 & 0.93 & 0.93 & 0.000 \\
consecutive mean (EVS, RLT) & 0.990 & 88 & 1.00 & 1.00 & 0.00 & 0.05 & 1.00 & 1.00 & 0.000 \\
consecutive cosine (DyCoke, FrameFusion) & 0.990 & 88 & 0.67 & 0.57 & 0.05 & 0.05 & 0.65 & 0.33 & 0.000 \\
segment static (PruneVid) & \multicolumn{9}{l}{cannot reach the budget at this frame rate} \\
frame similarity (LongVU) & 0.990 & 88 & 0.00 & 0.00 & 0.05 & 0.40 & 0.01 & 0.07 & 0.014 \\
consecutive + refresh every 60 & 0.985 & 132 & 0.03 & 0.00 & 0.45 & 1.00 & 0.01 & 0.00 & 0.014 \\
uniform frames & 0.992 & 66 & 0.03 & 0.00 & 0.30 & 0.70 & 0.01 & 0.07 & 0.008 \\
\bottomrule
\end{tabular}
}
\caption{Recall of planted events at a matched budget, CAVIAR lobby at 25~fps (top) and VIRAT at 10~fps (bottom). Every rule is calibrated to the drop rate of certified quotient + multi-scale at $\Delta_0 = 80$ (to within 0.001; whole-frame rules at their nearest setting). ``Cannot reach the budget'' means no threshold of the rule drops that share of patches; PruneVid's segment rule tops out near 0.94. The 2~fps runs are in Table~\ref{tab:audit2fps}.}
\label{tab:audit}
\end{table}

\paragraph{Consecutive rules miss fades.} At the matched budget, the consecutive mean catches 0--15\% of fades on three of the four runs, as Proposition~\ref{prop:consecutive} predicts: each frame's increment of a 3-s fade to contrast 48 is under 1 grey level at 25~fps, far below any threshold that drops 98\% of patches. Its median $\rho$ on the fading patches is $0.000$, so the proposition's second clause, firing by noise, gives it nothing either. The consecutive cosine does no better on fades (0--25\%). Both catch some brief and moving events where the scene is clean: on VIRAT the consecutive mean catches every small one-frame object and 81--88\% of movers, because a one-frame object is a large consecutive difference, but only 3--17\% of the tiny $3{\times}3$ objects, which a mean over the patch dilutes. The exception that confirms the mechanism is VIRAT at 2~fps, where the consecutive mean catches every 3-s fade: at 2~fps a 3-s fade is six frames, and each step is 8 grey levels. A consecutive threshold is a threshold on speed, and it is blind exactly when the frame rate is high relative to the change.

\paragraph{Anchored rules see fades and can miss small events.} Rules that compare against an older anchor see fades. Uniform sampling catches every fade on the lobby at 25~fps, and the consecutive rule with a refresh every 60 frames catches 75\% of 3-s fades and all 12-s fades there. But both miss almost every small, tiny and moving event (0--13\%): a periodic anchor sees what persists, not what is brief. LongVU's window similarity compares a frame only with the other frames of its 8-frame window, which is short against a fade, and it judges whole frames: it catches 0--15\% of fades and 0--7\% of small, tiny and moving events on all four runs, the window case of Proposition~\ref{prop:consecutive}. PruneVid's segment rule cannot reach the budget at all (it tops out near 0.94); at its tightest setting it catches every slow fade and 7--10\% of small events. The anchored mean, the anchored rule with no certificate, catches every fade on all four runs, as an anchor must, but its statistic dilutes small objects: it catches 3--20\% of tiny objects, 41--83\% of movers and, on the lobby, 27--30\% of small ones; a mean over 256 pixels dilutes a $5{\times}5$ object by a factor of ten.

\paragraph{The certified rule sees both.} Quotient + multi-scale catches 100\% of small, tiny and 3-s fade events, 90--95\% of 12-s fades and 90--93\% of movers on all four runs, at the same budget. It misses faint persistent objects of contrast 12 (0--20\%), which is below its single-frame certificate by design; the sequential layer, which only adds keeps, catches all of them (100\% on all four runs) for at most 0.1 points of drop rate. The range rule alone at the same budget is weak (0--30\% on every event): on real cameras sub-pixel edge jitter consumes its budget, which is why the quotient and multi-scale refinements exist. The certificate is not what produces the recall here; anchoring and a statistic that does not dilute small objects do. The certificate is what lets the recall be stated as a guarantee rather than measured on one set of planted events.
